% Titolo: Formazione del primo diario di bordo ed il relativo template
\documentclass[a4paper, 11pt]{article}

% --- PACCHETTI ---
\usepackage[utf8]{inputenc}
\usepackage[italian]{babel}
\usepackage[margin=2.5cm]{geometry}
\usepackage{graphicx}
\usepackage{tabularx}
\usepackage{hyperref}
\usepackage{fancyhdr}
\usepackage{array}
\usepackage{titlesec}
\usepackage{xcolor}

% --- CONFIGURAZIONE HEADER E FOOTER ---
\pagestyle{fancy}
\fancyhf{}
\fancyhead[L]{Bit Crusher}
\fancyhead[R]{Verbale Interno N. 3} 
\fancyfoot[C]{\thepage}
\renewcommand{\headrulewidth}{0.4pt}

% Configurazione link indice
\hypersetup{
    colorlinks=true,
    linkcolor=black,
    urlcolor=blue
}

\begin{document}

% ==========================================
% FRONTESPIZIO
% ==========================================
\begin{titlepage}
    \centering
    
    % --- LOGHI ---
    % Sostituisci i \framebox con i comandi \includegraphics corretti
    \begin{minipage}{0.45\textwidth}
        \centering
        % Esempio: \includegraphics[width=0.7\textwidth]{logo_unipd.png}
        \includegraphics[width=0.7\textwidth]{unipd_logo.png}
    \end{minipage}
    \hfill
    \begin{minipage}{0.45\textwidth}
        \centering
        % Esempio: \includegraphics[width=0.7\textwidth]{WhatsApp Image 2026-10-15 at 17.29.12.jpeg}
        \includegraphics[width=0.7\textwidth]{Logo-BitCrusher-Nero.png}
    \end{minipage}

    \vspace{1.5cm}
    {\Large \textbf{UNIVERSITÀ DEGLI STUDI DI PADOVA}} \\
    \vspace{0.3cm}
    {\large Corso di Laurea in Informatica} \\
    \vspace{0.3cm}
    {\large Insegnamento di Ingegneria del Software} \\
    \vspace{0.3cm}
    {\large Anno Accademico 2026/2027} \\

    \vspace{2cm}
    {\Huge \textbf{Verbale di Riunione}} \\
    \vspace{0.5cm}
    {\Large Riunione Interna \#3}

    \vspace{2cm}
    \begin{flushleft}
        \large
        \textbf{Gruppo:} Bit Crusher \\
        \textbf{Email:} bitcrusherteam@gmail.com \\
        \textbf{Versione:} v1.0.0 \\
        \textbf{Data della riunione:} 2026-10-24 \\
        \textbf{Data ultima modifica:} 24 ottobre 2026
    \end{flushleft}
    \vfill
\end{titlepage}

\newpage

% ==========================================
% REGISTRO DELLE MODIFICHE
% ==========================================
\section*{Registro delle Modifiche}
\vspace{0.5cm}
\begin{center}
\renewcommand{\arraystretch}{1.5}
\begin{tabularx}{\textwidth}{|l|l|>{\raggedright\arraybackslash}p{3.5cm}|>{\raggedright\arraybackslash}p{3.5cm}|X|}
%modificato per permettere ai campi nella tabella di andare a capo
\hline
\textbf{Versione} & \textbf{Data} & \textbf{Autore} & \textbf{Verificatore} & \textbf{Descrizione} \\ \hline
1.0.0 & 2026-11-03 & Alex Singh & - & Approvazione per il rilascio. \\ \hline
0.2.0 & 2026-11-03 & Antonio Tufano & Riccardo Pozzato & Aggiornamento logo UniPD e adeguamento registro modifiche. \\ \hline
0.1.0 & 2026-10-24 & Pietro Leonardo Acampora & Marco Bollettin, Antonio Tufano & Stesura verbale di riunione \\ \hline

\end{tabularx}
\end{center}

\newpage

% ==========================================
% INDICE
% ==========================================
\tableofcontents
\newpage
% ==========================================
% 1. INFORMAZIONI GENERALI
% ==========================================
\section{Informazioni Generali}
\begin{itemize}
    \item \textbf{Luogo/Piattaforma:} Riunione Discord
    \item \textbf{Data:} 2026-10-24
    \item \textbf{Orario di inizio:} 17:30
    \item \textbf{Orario di fine:} 20:15
    \item \textbf{Responsabile:} Riccardo Pozzato
    \item \textbf{Scriba:} Riccardo Pozzato
\end{itemize}

\vspace{0.5cm}
\textbf{Partecipanti presenti:} \\
\vspace{0.2cm}
\renewcommand{\arraystretch}{1.2}
\begin{tabularx}{0.8\textwidth}{|X|c|}
\hline
\textbf{Nome e Cognome} & \textbf{Matricola} \\ \hline
Pietro Leonardo Acampora & 2146762 \\ \hline
Marco Bollettin & 2147955 \\ \hline
Alexandru Mitu & 2101083 \\ \hline
Riccardo Pozzato & 2147988 \\ \hline
Alex Singh & 2148621 \\ \hline
Antonio Tufano & 2137972 \\ \hline
Erik Giuliano & 2138012 \\ \hline
\end{tabularx}

\vspace{0.5cm}
\textbf{Partecipanti assenti:} \\
\vspace{0.2cm}
\renewcommand{\arraystretch}{1.2}
\begin{tabularx}{0.8\textwidth}{|X|c|}
\hline
\textbf{Nome e Cognome} & \textbf{Matricola} \\ \hline
 Nessuno & - \\ \hline
\end{tabularx}

% ==========================================
% 2. ORDINE DEL GIORNO
% ==========================================
\vspace{0.5cm}
\section{Ordine del Giorno}
Gli argomenti previsti per la riunione odierna sono:
\begin{enumerate}
    \item Creazione del template .pptx per i diari di bordo
    \item Creazione diario di bordo N.1
    \item Stesura presentazione diario di bordo
    \item Scelta rappresentante per presentazione diario di bordo N.1
    \item Discussione riguardo le risposte ricevute dalle aziende in merito alle domande fatte
\end{enumerate}

% ==========================================
% 3. RESOCONTO DELLA RIUNIONE
% ==========================================
\newpage
\section{Resoconto della Riunione}

\subsection{Verifica delle azioni precedenti}
\begin{itemize}
    \item \textbf{VI-1.1 del verbale 02} Il gruppo ha esaminato le risposte delle aziende ai quesiti inviati. Il dialogo è proseguito con il primo incontro con VarGroup, documentato nel verbale esterno inserito il 27/10/2026 (\href{https://github.com/BitCrusher-Team/Progetto_SWE_26-27/pull/14}{PR \#14}).
\end{itemize}

\subsection{Creazione del template .pptx per i diari di bordo}
Il gruppo ha discusso e definito la struttura, lo stile e la grafica del template in formato PowerPoint da utilizzare per le presentazioni dei futuri diari di bordo.

\subsection{Creazione diario di bordo N.1}
È stato redatto in modo collaborativo il contenuto del primo diario di bordo, riassumendo i risultati raggiunti, le decisioni prese e i problemi riscontrati nel periodo corrente.

\subsection{Stesura presentazione diario di bordo}
Sono state preparate e impaginate le slide per la presentazione del diario di bordo N.1, inserendo i contenuti discussi all'interno del nuovo template.

\subsection{Scelta rappresentante per presentazione diario di bordo N.1}
I membri del gruppo si sono confrontati e hanno selezionato formalmente la persona incaricata di esporre la presentazione durante la prima revisione di avanzamento.

\subsection{Discussione riguardo le risposte ricevute dalle aziende in merito alle domande fatte}
Sono state lette e analizzate congiuntamente le email di risposta fornite dalle aziende proponenti, valutando i chiarimenti ottenuti in merito ai dubbi relativi ai capitolati.

\subsection{Decisioni prese}
\begin{itemize}
    \item \textbf{V-3.1} Scelta la tecnologia da usare per la redazione e presentazione dei diari di bordo
    \item \textbf{V-3.2} Definite le convenzioni per l'approccio trunk-based development sul repository
\end{itemize}

% ==========================================
% 4. AZIONI DA SVOLGERE (TO-DO)
% ==========================================
\vspace{0.5cm}
\section{Azioni da Svolgere (To-Do)}
Nessun task assegnato al termine di questa riunione. Le attività elencate sono ricostruite dalle modifiche successive al repository.

\begin{tabularx}{\textwidth}{|c|X|l|l|}
\hline
\textbf{Codice Rif.} & \textbf{Descrizione Task} & \textbf{Assegnatario} & \textbf{Scadenza} \\ \hline
--- & Pubblicare nel repository il primo diario di bordo e il template PowerPoint. Entrambi sono stati aggiunti il 30/10/2026 (\href{https://github.com/BitCrusher-Team/Progetto_SWE_26-27/pull/26}{PR \#26}). & --- & --- \\ \hline
--- & Rendere operative le convenzioni per il lavoro a stack. Il 27/10/2026 sono stati aggiunti gli alias e lo script per gestire gli stack con \texttt{git spr} (\href{https://github.com/BitCrusher-Team/Progetto_SWE_26-27/commit/45c68ee}{commit \texttt{45c68ee}}). & --- & --- \\ \hline
\end{tabularx}

% ==========================================
% 5. PROSSIMA RIUNIONE
% ==========================================
\vspace{0.5cm}
\section{Prossima Riunione}
\begin{itemize}
    \item \textbf{Data:} Da concordare
    \item \textbf{Ora:} Da concordare
    \item \textbf{OdG preliminare:} Da concordare
\end{itemize}

\end{document}